\section{问题二：无扰流正常工况校准}

\subsection{问题分析}

附件中标记为无扰流参考工况的窗口反映了仪表在正常速度剖面下随流量变化的固有偏差。因此，该部分数据可以作为正常工况校准依据。

\subsection{问题建模}

令问题一中选定的 Phys6 工程积分估计量为
\begin{equation}
V_{b,i}=V_{\mathrm{Phys6},i}.
\end{equation}
在无扰流状态下，可引入只依赖流量规模的百分数正常校准曲线 $c_0^{(\%)}(q)$，得到
\begin{equation}
V_{N,i}=V_{b,i}\left[1+\frac{c_0^{(\%)}(q_i)}{100}\right],
\label{eq:d0_curve}
\end{equation}
其中 $q_i$ 为归一化流量点。程序中采用的归一化公式为
\begin{equation}
q_i=\frac{\mathrm{flow\ point}_i-50}{30}.
\end{equation}
本文采用二次平滑形式
\begin{equation}
c_0^{(\%)}(q)=b_0+b_1q+b_2q^2,
\end{equation}
并在 D0 窗口上用带岭惩罚的最小二乘估计参数，以避免少量无扰流样本造成过拟合。

这里的 $q_i$ 取题目提供的名义流量点并归一化，表示运行时流量规模，不读取标准装置体积。

在标签盲示例划分中，训练集含有 $7$ 个 D0 窗口。表中系数只由这些训练窗口估计，$41$ 个验证窗口不参与参数求解。结果见表~\ref{tab:d0_curve_coef}。这里的曲线直接修正 Phys6 基线，与问题四在 KRR 之后拟合的 D0 残余后校准曲线不同。

\begin{table}[H]
\centering
\caption{标签盲示例划分训练集上的 Phys6 无扰流校准曲线系数}
\label{tab:d0_curve_coef}
\begin{tabular}{lcc}
\toprule
项 & 系数(\%) & 曲线形式 \\
\midrule
$b_0$ & 0.402727 & $c_0^{(\%)}(q)=b_0+b_1q+b_2q^2$ \\
$b_1$ & -0.033031 & 同上 \\
$b_2$ & -0.007535 & 同上 \\
\bottomrule
\end{tabular}
\end{table}

\subsection{结果分析}

无扰流校准曲线只描述正常流态中随流量变化的系统偏差，不能表达弯管、旋流或非充分发展流造成的附加剖面畸变。因此，本节重点考察正常工况线性度与流量点残差，不能把该曲线直接视为扰流补偿模型。

在标签盲示例划分消融中，仅使用 Phys6 时
\begin{equation}
u_{\mathrm{nor},L}=0.462931\% .
\end{equation}
加入 D0 曲线后，
\begin{equation}
u_{\mathrm{nor},L}=0.055310\% ,
\end{equation}
说明正常流量点系统偏差被显著削弱。但同一消融中，
\begin{equation}
\mathrm{MAE}:0.179991\%\rightarrow 0.260279\%,
\qquad
u_{\mathrm{nor},d}:0.217668\%\rightarrow 0.218491\%,
\end{equation}
没有同步改善。这说明无扰流曲线能够削弱正常流量点的系统偏差，但不能代表扰流状态下的补偿效果。

无扰流内部实验还表明，带对称权重约束的 OWICS/Gauss-Jacobi 先验与剖面修正模型在正常窗口上误差较低；但直接作用于其余扰流窗口后，误差明显增大。因此，正常工况下得到的最优积分或校准关系不能直接外推到畸变剖面。

综上，问题二得到一条只描述正常工况流量相关偏差的校准曲线，为下一步分析扰流造成的额外剖面变化提供了一个基准。

从 D0 内部逐流量点结果看，校准模型在 $40,50,60,70,80$ 五个流量点的 MAE 分别为 $0.046174\%$、$0.035041\%$、$0.027378\%$、$0.057486\%$ 和 $0.036046\%$。其中 $70$ 流量点残余误差最大，仍需更多正常锚点；$40$ 流量点仅有一个窗口，不能据此判断重复性。由于 D0 全部来自同一天，上述结果属于正常工况内部拟合分析，而不是跨日期验证。
